\section{逮捕}
\begin{enumerate}
  \item 法检决定公安执行。
  \item 公安报检察院 \yjt 批准速捕，检察院自己想捕 \yjt 决定逮捕。
  \item 适用条件。基本条件：
        \begin{enumerate}
          \item 有证据。发生了，他干了，已有查证属实了。
          \item 有徒刑。
          \item 有危险。新罪逃杀串毁证，国公秩序有危险，还要打击报复他。
          \item 径行逮捕（直接）。10 年以上。太严重或太危险。
          \item 转化型逮捕。违反取保候审、监视居住规定，情节严重的。
        \end{enumerate}
  \item \zd{检察院批捕程序}。
        \begin{enumerate}
          \item 报请批捕。公安提请批准逮捕书报同级检察院。
          \item 决定期限。已拘留 7 天，没有 15 天；重大、复杂不过 20 天。
          \item 结果。批捕 \hz 不批捕。
          \item 不批捕复议复核。是原条件（比如补充了新证据，条件就变了）不变的情况下。
          \item 应当捕前讯问。\zd{疑}问、当\zd{面}陈述、重大\zd{违}法、重大疑\zd{难、认罪认罚；小孩聋哑半疯傻}。
          \item 听取意见。未成年人，应当听取辩护律师意见。
          \item 排非程序。
        \end{enumerate}
  \item 检察院决定逮捕。公安监察移送审查起诉的案件；自侦案件。
  \item 法院决定逮捕。公诉案件；自诉案件。
  \item 执行程序。
        \begin{enumerate}
          \item 逮捕证没有例外；2 人以上；\zd{立即}送看；24 小时内通知家属、应当询问（谁想捕，谁通知、讯问）。
        \end{enumerate}
  \item 特殊主体。
        \begin{enumerate}
          \item 人大代表的强制措施适用报请人大常委会 \hz 主席团的许可之后。
          \item 外国人。\zd{危}害国家安全、影响国家政治外\zd{交}、法律适用有疑\zd{难找最高检，其他先批后备案}。
          \item 危害国家安全。报上一级备案。
        \end{enumerate}
  \item \zd{羁押必要性审查}。\zd{被逮捕}的犯罪嫌疑人、被告人。
        \begin{enumerate}
          \item \zd{捕诉部门}；原则公开审查，三密除外。
          \item 达不到逮捕条件，应当变更逮捕。\zd{无证据、无徒刑、务必要、已超期}。
        \end{enumerate}
  \item 法院应当变更。\zd{无罪无责不处罚；量刑非羁押；刑期已折抵；办案已超期}。
  \item 自动解除。\zd{单出附加生效后；主刑开始执行后}。
\end{enumerate}
